Transit bus operators face competing demands driven by financial, customer, and regulatory concerns. As subsidized public infrastructure transit agencies often operate at a loss, only directly recovering a portion of their expenses as revenue \cite{Litman_2014_JPT, Taylor_2009_JPL}, a situation worsened by the COVID-19 pandemic \cite{Siddiq_2023_PWMP}. American transit ridership has been in decline for much of the 21\textsuperscript{st} century \cite{Wasserman_2022_JPT} due to competition from other modes \cite{Erhardt_2022_TRA, Pan_2022_POM}. Additional stress is provided by regulatory requirements to reduce fleet emissions \cite{Jeffers_2022_OSTI, CARB_2018_CARB, FTA_2022_FTA} by adopting \glspl{zev}. Provision of affordable, clean, and desirable public transit is a difficult and multifaceted task requiring careful planning.

In strategic planning transit operators must make several related decisions. These are (1) what routes to offer, (2) what timings and frequencies to offer for these routes, and (3) what vehicles and other facilities and equipment to purchase. These decisions will not necessarily be made around the same objectives. Decisions (1) and (2) are conveniently formulated as optimization problems and serve as standard problems in the optimization research field. However, transit agencies often determine their routes based on unconventional and even un-quantifiable factors, seeking to serve diverse customer priorities \cite{Mers_2024_TRR}. Because transit-riders value predictability and consistency \cite{Watkins_2011_TRA, Berrebi_2021_TRA}, routes and schedules are static and changes are introduced slowly. By contrast, changes to fleets are driven by funding cycles, regulatory changes, and scheduled equipment turnover and, thus, may happen quickly and out-of-step with changes to routes and schedules. For this reason, procurement of vehicles and supporting facilities and equipment is considered independent of route design.

The resulting problem is the \gls{vsfmp}. The \gls{vsfmp} encompasses optimization problems in which a set of location and time specified demand nodes must be serviced by a fleet of vehicles of unknown composition operating from a set of depots. The \gls{vsfmp} is a sub-problem of the \gls{vsp} \cite{Bunte_2009_PT} and the related \gls{vrp} \cite{Konstantakopoulos_2022_OR, Sabet_2022_IA}, both known to be NP-hard \glspl{cop}. The \gls{vsfmp} starts with a set of trips defined by start and finish locations and times and route information such as energy consumption and expected ridership. If these trips are relatively short, then substantial savings can be made by serving multiple trips consecutively with the same vehicle before returning to depot. Chains of consecutive trips serviced by one vehicle before returning to depot are called tours. Finding the best set of tours is the core of the \gls{vsp}. Where energy-resupply time and supply equipment costs are non-trivial, a second assignment problem must be solved related to supply events and supply equipment. For the \gls{vsfmp} the number of vehicles, vehicle type mix, number of supply equipment, supply equipment mix, and allocation of vehicles and supply equipment to depots are all unknowns. When formulated as a \gls{milp}, the \gls{vsfmp} is subject to extreme combinatorial explosion. Run-time scales exponentially with the size of the bus network as well as the number of different technologies being considered. The run-time complexity of \gls{milp} formulations of the \gls{vsfmp} reduce the practical utility of existing methodology for medium and large bus operators.

It is worth discussing the reasons that the \gls{vsp} and \gls{vrp} are so difficult to optimize. A generic formulation of the \glspl{vsfmp} is provided in \cite{Rogge_2018_AE} which is built around several assignment matrices. For each vehicle type under consideration, an assignment matrix $\mathcal{A}$ connects predecessor and successor trips. For a bus network defined by directed graph $G = \{V \cup D, E\}$ where $V$ is the set of trips, $D$ is the set of depots, and $E$ is the set of feasible edges, $\mathcal{A}$ is necessarily $|V| + |D|$ by $|V| + |D|$ and contains, at minimum, $|E|$ binary variables. Since trips may have only one predecessor and only one successor each non-depot row and column must sum to $1$. This constraint introduces substantial non-convexity to the problem. When energy supply is considered, an additional assignment of supply ports to supply events must be made for all supply port types with similar unity constraints. The number of binary variables scales with $(|V| + |D|)^2$ and linearly with the number of vehicle and port types considered.

For highly non-convex integer optimization, \gls{ga} provides several key advantages compared to optimal solver methods \cite{Hajela_1990_AIAA}. These advantages stem from \gls{ga} not relying on a gradient search which allows for more complete design exploration and minimizes the risk of settling in local minima. From a computational perspective, the fixed population size means that \gls{ga} solutions may require far fewer function evaluations than traditional optimization methodology to arrive at a near-optimal solution. The downside is that it is not possible to compute an optimality gap for \gls{ga} and thus not possible to know if the solution arrived at is near-optimal. Additionally, as \gls{ga} relies on pseudo-random number generation, there is the possibility of initial seed influence determining the final solution. Increased population size, mutation, and population-migration strategies mitigate this issue to some degree. Another key advantage of \gls{ga} is that it can be extended to multi-objective (Pareto) optimization. The most common multi-objective implementation of \gls{ga} is \gls{nsga} \cite{Deb_2002_TEC}. \gls{nsga} retains most of the single-objective \gls{ga} method but, rather than computing fitness along a single objective, fitness is determined by Pareto dominance rank and crowding distance within a given rank. \gls{nsga} is particularly useful as it allows for intractable, non-convex, multi-objective \glspl{cop} to be readily solved in a reasonable amount of time for realistically scaled systems.

This paper presents a novel and computationally efficient decision support tool that advances transit operator strategic planning by performing efficient Pareto optimization over transit systems defined by routes, trips, and depots. By revealing a multi-objective Pareto front rather than a single optimization point, the tool allows decision-makers to make fully informed decisions. This method works with any number of vehicle types, any number of supply infrastructure types, any number of depots greater than or equal to one, and any number of objective functions. Further, run-time is almost entirely determined by network size and \gls{ga} hyper-parameters and is not meaningfully impacted by the number of vehicle types, infrastructure types, depots, or objectives. Implementations of the method are provided in Python (flexible implementation) and Rust (efficient implementation). This paper surveys the existing literature, describes the optimization method providing background on the \gls{nsga} on which it is based and where it diverges, and provides an illustrative example based on the UC Davis Unitrans bus network, which is a medium sized transit bus network.

\section{Literature Review}

Traditionally, energy resupply requirements have not been central to operational problems such as the \gls{vsp} and \gls{vrp}. The introduction of \glspl{bev} and \glspl{fcev}, with non-negligible energy resupply times and limited resupply locations, necessitated serious consideration of energy resupply in operations. This resulted in added problem complexity including non-convex constraints. One approach is to omit the vehicle scheduling entirely. On the theory that vehicle scheduling is either fixed or optimized around unrelated objectives, the procurement and allocation of supply equipment can be tractably solved for as an \gls{ap} as in \cite{Foda_2023_E, Foda_2024_TP, Gkiotsalitis_2025_AE}. In order to render the full optimization tractable, \cite{Wang_2022_TRD} introduced a linearized optimal formulation based on a set of restricting assumptions. With this formulation, a bus network of 4 lines and 1 depot resulted in a model with nearly 60,000 binary variables. Using Gurobi and a reasonable desktop computer, a solution with an optimality gap of less than 4\% was obtained after 30 minutes. This result is simultaneously impressive and demonstrative of the computational intractability of the problem at realistic scales.

Many researchers have chosen to deal with problem complexity via the use of meta-heuristics, specifically variants of \gls{ga}. Operations problems such as the \gls{vrp} and the related \gls{sbrp} are easily encoded for \gls{ga}. For these problems, demand nodes can be grouped into subsets and serviced in an optimal order as determined by a local search heuristic \cite{Jozefowiez_2006_AE, Sooktip_2015_CC, Hou_2022_EL}. This holds for the \gls{vrptw} where failure to adhere to delivery windows results in higher costs rather than infeasibility \cite{Mendes_2016_ITSC, Tang_2021_ITITS, Feng_2024_ICECAI, Wang_2020_ESA}. Encoding becomes more difficult for \gls{vsp} type problems where demand nodes must be serviced at specific times. In this case, many groupings of demand nodes will be infeasible.

Several approaches are seen in literature. One approach is to encode the \gls{milp} directly as in \cite{Li_2020_JAT}. This results in a large number of variables and a large number of infeasible offspring in each generation. These issues are somewhat mitigated with an adaptive crossover scheme. However, \cite{Li_2020_JAT} notes that the algorithm is not particularly computationally efficient and scales poorly. A second approach is the "bin-packing" method used in \cite{Yao_2020_SCS} in which chromosomes represent a random ordering of service trips. Service trips are then assigned to buses based on the ordering using a greedy heuristic. Where trips cannot be fit into the schedule of an existing bus, a new bus is created. Although assignment is fully determined by the chromosome the influence of the chromosome is highly indirect. Given the nature of the assignment heuristic, the earlier parts of the chromosome have much more leverage over chromosome fitness than later parts. A third approach is proposed in \cite{Rogge_2018_AE} which encodes a grouping of trips into feasible tours and then solves the resulting assignment problem as a \gls{milp} using an optimal solver. The primary purpose of the grouping in \cite{Rogge_2018_AE} is to reduce the size of the \gls{milp} compared to a trip-assignment problem as in \cite{Wang_2022_TRD}. A trip insertion heuristic is adopted for initial generation, crossover, and mutation of the tours. Crossover and mutation result in unassigned trips which the insertion heuristic inserts in the lowest-cost feasible slot. The insertion heuristic guarantees that there will never be an infeasible solution but tends to over-determine population evolution. As a result, \cite{Rogge_2018_AE} utilizes a population-migration version of \gls{ga} to preserve diversity and avoid premature convergence. Although more efficient than a traditional \gls{milp} solution, the method in \cite{Rogge_2018_AE} does not scale well enough for use with very large networks. The cited works constitute meaningful contributions to the \gls{vsfmp} literature. However, they do not provide a straightforward generic implementation, and they address only single objective formulations.

As bus fleet operators are concerned with multiple objectives simultaneously, true multi-objective optimization methods are of higher utility than single-objective methods. Pareto optimization methodology (methods which reveal the Pareto front rather than a single optimum) has been applied to operations research for decades \cite{Baita_2000_CER}. The most common method for Pareto optimization is \gls{nsga} \cite{Deb_2002_TEC} which has been applied extensively in the operations domain \cite{Verma_2021_IA}. Although \gls{nsga} has been applied to the \gls{vrp}, the additional complexities of encoding the \gls{vsp} for \gls{ga} have prevented a generic implementation Pareto optimization method for the \gls{vsfmp}.

This paper represents a significant advancement of the state-of-the-art by providing a generic implementation of \gls{nsga} for the \gls{vsfmp}. This implementation is generic in that it does not rely on any specific objective function or constraint. The implementation is also highly efficient, achieving rapid convergence while maintaining or extending field spread. The efficiency of the method is demonstrated in the Results section and can be verified using the provided Python library. Finally, the provided implementation has favorable scaling properties and can be used to optimize realistically scaled bus networks in a reasonable amount of time.